\scanpage{13}

% archive sheet number: 12

rendant commutatifs les diagrammes

\diagram{\begin{tikzcd}[ampersand replacement=\&]
\text{(i)}\quad C\arrow[r,"q_C"]\arrow[rr,bend right=14,"p_1"']\& B\otimes C\arrow[r,"{(g_C,p_2)}"]\& C\otimes C
\end{tikzcd}}

\deleted{\unsure{(ii)}}

\diagram{\begin{tikzcd}[ampersand replacement=\&]
\text{(ii)}\quad B\arrow[r,"g_C"]\arrow[rr,bend right=20,"p_1"']\& C\otimes C\arrow[r,"{\varphi=(q_C,p_2)}"]\& B\otimes C
\end{tikzcd}}

D'autre part, les données \unsure{$(*)$|$*)$} définissent

\[
\psi:\operatorname{Hom}_A(B,C)\longrightarrow\operatorname{Hom}_A(C,C)
\]

en posant $\psi(v)=\Delta_C(v\otimes\mathrm{id}_C)q_C$ :

\diagram{\begin{tikzcd}[ampersand replacement=\&]
C\arrow[r,"q_C"]\arrow[rrr,bend right=12,"{\psi(v)}"']\& B\otimes C\arrow[r,"{v\otimes\mathrm{id}_C}"]\& \deleted{\unsure{B}}C\otimes C\arrow[r,"{\Delta_C}"]\&C
\end{tikzcd}}

Je dis que \uline{si} $\varphi$ \uline{est un isomorphisme}, $\psi$ \uline{en est un}.

On va en effet construire, utilisant $g_C$, un hom

\[
\psi':\operatorname{Hom}(C,C)\longrightarrow\operatorname{Hom}(B,C)
\]

qui sera inverse de $\psi$. On pose

\[
\psi'(w)=\Delta_C(w\otimes\mathrm{id}_C)g_C\quad :
\]

\diagram{\begin{tikzcd}[ampersand replacement=\&]
B\arrow[r,"g_C"]\arrow[rrr,bend right=13,"{\psi'(w)}"']\& C\otimes C\arrow[r,"{w\otimes\mathrm{id}_C}"]\& C\otimes C\arrow[r,"{\Delta_C}"]\& C
\end{tikzcd}}

On doit prouver $\psi'(\psi(\unsure{w|v}))=v$, $\psi\psi'(w)=w$. C'est facile :

\[
\deleted{\psi\psi'(w)=\Delta_{\unsure{C}}[(\Delta_C(w\otimes\mathrm{id}_{\unsure{C}})g_C)\otimes\mathrm{id}_C]q_C}
\]
